\chapter{Symmetry, Groups, and Lie Algebras}

\section{What a group is}

A group is a set with a multiplication rule satisfying four conditions:
\[
  \text{closure},\quad \text{associativity},\quad \text{identity},\quad \text{inverse}.
\]
Rotations of the plane form a group.  A rotation by angle \(\theta\) is
\[
  R(\theta)=
  \begin{pmatrix}
  \cos\theta&-\sin\theta\\
  \sin\theta&\cos\theta
  \end{pmatrix}.
\]
Multiplying rotations adds angles:
\[
  R(\theta)R(\phi)=R(\theta+\phi).
\]
The identity is \(R(0)\), and the inverse is \(R(-\theta)\).

\section{Continuous symmetries and generators}

For small \(\theta\),
\[
  R(\theta)
  =
  I+\theta
  \begin{pmatrix}
  0&-1\\
  1&0
  \end{pmatrix}
  +O(\theta^2).
\]
In quantum notation one often writes
\[
  U(\theta)=e^{-\ii\theta T},
\]
where \(T\) is the generator.  Expanding,
\[
  U(\theta)=1-\ii\theta T+O(\theta^2).
\]
The generator contains the infinitesimal action; the exponential builds finite transformations.

\section{U(1)}

The group \(U(1)\) consists of phases
\[
  e^{\ii\alpha}.
\]
It acts on a complex scalar by
\[
  \Phi\to e^{\ii q\alpha}\Phi.
\]
The number \(q\) is the charge in units of the chosen coupling.  The infinitesimal form is
\[
  \delta\Phi=\ii q\alpha\Phi.
\]
Global \(U(1)\) symmetry produced the conserved current of the complex scalar field.  Making \(\alpha\) depend on spacetime will force us to introduce electromagnetism.

\section{SU(2) and Pauli matrices}

The group \(SU(2)\) consists of \(2\times2\) unitary matrices with determinant \(1\).  Its generators can be chosen as
\[
  T^a=\frac{\sigma^a}{2},
\]
where
\[
  \sigma^1=
  \begin{pmatrix}0&1\\1&0\end{pmatrix},
  \quad
  \sigma^2=
  \begin{pmatrix}0&-\ii\\ \ii&0\end{pmatrix},
  \quad
  \sigma^3=
  \begin{pmatrix}1&0\\0&-1\end{pmatrix}.
\]
Direct multiplication gives
\[
  [\sigma^a,\sigma^b]=2\ii\epsilon^{abc}\sigma^c.
\]
Therefore
\[
  [T^a,T^b]=\ii\epsilon^{abc}T^c.
\]
The constants \(\epsilon^{abc}\) are structure constants.  They encode the Lie algebra.

\section{SU(N)}

For a general non-Abelian group with generators \(T^a\),
\[
  [T^a,T^b]=\ii f^{abc}T^c.
\]
If all \(f^{abc}=0\), the group is Abelian.  If not, transformations fail to commute.  This failure is the source of self-interactions in Yang--Mills theory.

A field may transform in a representation:
\[
  \psi\to U\psi,
  \qquad
  U=e^{-\ii\alpha^aT^a}.
\]
The matrices \(T^a\) depend on the representation.  The algebraic relation above must hold in each representation.

\section{Noether currents from internal symmetries}

Let fields \(\phi_i\) transform as
\[
  \delta\phi_i=\epsilon\,\Delta_i(\phi).
\]
If the Lagrangian is invariant, the Noether current is
\[
  j^\mu=
  \sum_i
  \frac{\p\Lag}{\p(\p_\mu\phi_i)}
  \Delta_i.
\]
For a complex scalar with \(U(1)\) transformation
\[
  \delta\Phi=\ii\epsilon\Phi,
  \qquad
  \delta\Phi^*=-\ii\epsilon\Phi^*,
\]
and
\[
  \Lag=\p_\mu\Phi^*\p^\mu\Phi-m^2\Phi^*\Phi,
\]
we get
\[
  j^\mu=
  \frac{\p\Lag}{\p(\p_\mu\Phi)}\ii\Phi
  +
  \frac{\p\Lag}{\p(\p_\mu\Phi^*)}(-\ii\Phi^*)
  =
  \ii\left(\p^\mu\Phi^*\Phi-\p^\mu\Phi\,\Phi^*\right).
\]
This differs by an overall sign from the earlier convention depending on whether the infinitesimal parameter is assigned to \(\Phi\) or \(\Phi^*\).  The conserved charge changes sign with the convention, not the physics.

\section{Local symmetry as a demand}

Suppose a Lagrangian is invariant under constant \(U=e^{-\ii\alpha^aT^a}\).  If \(\alpha^a=\alpha^a(x)\), then
\[
  \p_\mu(U\psi)=(\p_\mu U)\psi+U\p_\mu\psi.
\]
The extra \((\p_\mu U)\psi\) spoils the simple transformation law.  Gauge theory introduces a covariant derivative \(D_\mu\) such that
\[
  D_\mu\psi\to U D_\mu\psi.
\]
Everything in gauge theory follows from building this derivative and its curvature.

\section*{Exercises}
\addcontentsline{toc}{section}{Exercises}

\begin{exercise}
Verify that \(\det R(\theta)=1\) for the two-dimensional rotation matrix.
\begin{answer}
\(\cos^2\theta+\sin^2\theta=1\).
\end{answer}
\end{exercise}

\begin{exercise}
Compute \([\sigma^1,\sigma^2]\).
\begin{answer}
\(\sigma^1\sigma^2=\ii\sigma^3\), \(\sigma^2\sigma^1=-\ii\sigma^3\), so the commutator is \(2\ii\sigma^3\).
\end{answer}
\end{exercise}
